\begin{proof}[Proof of \cref{obj:lem:shrinking-cone-dual}]
\begin{enumerate}
\item Define the reciprocal target and its slope parameter by
\[
c_\epsilon:=1-\epsilon,
\qquad
g_{\mathrm{rec}}(v):=\frac{\alpha}{\alpha+c_\epsilon v}
\quad (v\in[0,B]).
\]
The hypotheses give
\[
\alpha>0,
\qquad
c_\epsilon\ge\frac34>0,
\qquad
\alpha+c_\epsilon v\ge\alpha>0
\quad (v\in[0,B]).
\]
Thus \(g_{\mathrm{rec}}\) is continuous on \([0,B]\). % lean: CausalSmith.Stat.AnnotationRarearmFrontier.cone_reciprocal_continuous

\item We first select a polynomial attaining the best uniform approximation error. For a continuous real function \(\varphi\) on a nonempty compact interval, write
\[
\|\varphi\|_{\infty,[a_-,a_+]}
:=
\sup_{v\in[a_-,a_+]}|\varphi(v)|.
\]
Define
\[
\mathcal V_{\mathrm{app}}
:=
\left\{
Q|_{[0,B]}:
Q\text{ is a real polynomial of degree at most }L
\right\},
\]
\[
\rho_{\mathrm{app}}
:=
2\|g_{\mathrm{rec}}\|_{\infty,[0,B]}+1,
\qquad
\mathcal K_{\mathrm{app}}
:=
\left\{
\varphi\in\mathcal V_{\mathrm{app}}:
\|\varphi\|_{\infty,[0,B]}\le\rho_{\mathrm{app}}
\right\}.
\]
The space \(\mathcal V_{\mathrm{app}}\) is finite dimensional, being spanned by the restrictions of \(1,v,\ldots,v^L\). Hence \(\mathcal K_{\mathrm{app}}\) is compact and contains zero. The continuous error functional therefore attains its minimum there at the restriction of a polynomial \(Q_{\mathrm{app}}\) of degree at most \(L\). Comparison with zero gives
\[
\|g_{\mathrm{rec}}-Q_{\mathrm{app}}\|_{\infty,[0,B]}
\le
\|g_{\mathrm{rec}}\|_{\infty,[0,B]}.
\]
For \(\varphi\in\mathcal K_{\mathrm{app}}\), minimality gives
\[
\|g_{\mathrm{rec}}-Q_{\mathrm{app}}\|_{\infty,[0,B]}
\le
\|g_{\mathrm{rec}}-\varphi\|_{\infty,[0,B]}.
\]
For \(\varphi\in\mathcal V_{\mathrm{app}}\setminus\mathcal K_{\mathrm{app}}\), the reverse triangle inequality gives
\[
\begin{aligned}
\|g_{\mathrm{rec}}-\varphi\|_{\infty,[0,B]}
&\ge
\|\varphi\|_{\infty,[0,B]}
-\|g_{\mathrm{rec}}\|_{\infty,[0,B]}\\
&>
\|g_{\mathrm{rec}}\|_{\infty,[0,B]}+1
\ge
\|g_{\mathrm{rec}}-Q_{\mathrm{app}}\|_{\infty,[0,B]}.
\end{aligned}
\]
Consequently, defining the best approximation error by
\[
E_{\mathrm{app}}
:=
\inf_{\substack{Q\text{ real polynomial}\\ \deg Q\le L}}
\|g_{\mathrm{rec}}-Q\|_{\infty,[0,B]},
\]
we have
\[
E_{\mathrm{app}}
=
\|g_{\mathrm{rec}}-Q_{\mathrm{app}}\|_{\infty,[0,B]}.
\]
% lean: Causalean.Mathlib.Analysis.FinitePolynomialAlternationDuality.exists_bestPolynomial

\item Suppose, for contradiction, that
\[
\|g_{\mathrm{rec}}-Q_{\mathrm{app}}\|_{\infty,[0,B]}<\frac18.
\]
Then, for every \(v\in[0,B]\),
\[
|g_{\mathrm{rec}}(v)-Q_{\mathrm{app}}(v)|<\frac18,
\qquad
0\le g_{\mathrm{rec}}(v)\le1.
\]
The triangle inequality therefore yields
\[
|Q_{\mathrm{app}}(v)|
\le
|Q_{\mathrm{app}}(v)-g_{\mathrm{rec}}(v)|
+g_{\mathrm{rec}}(v)
<\frac98,
\qquad
\|Q_{\mathrm{app}}\|_{\infty,[0,B]}\le\frac98.
\]
Define the comparison point by
\[
z_{\mathrm{rec}}:=\frac{\alpha}{c_\epsilon}.
\]
Since \(L^2\ge4\) and \(c_\epsilon\ge3/4\),
\[
0<z_{\mathrm{rec}}
=
\frac{B}{100L^2c_\epsilon}
\le\frac{B}{300}\le B.
\]
Direct evaluation gives
\[
g_{\mathrm{rec}}(0)=1,
\qquad
g_{\mathrm{rec}}(z_{\mathrm{rec}})=\frac12.
\]
The approximation bounds at these two points imply
\[
Q_{\mathrm{app}}(0)>\frac78,
\qquad
Q_{\mathrm{app}}(z_{\mathrm{rec}})<\frac58,
\qquad
Q_{\mathrm{app}}(0)-Q_{\mathrm{app}}(z_{\mathrm{rec}})>\frac14.
\]
By the mean value theorem, there exists \(v_{\mathrm{mvt}}\in(0,z_{\mathrm{rec}})\) such that
\[
Q_{\mathrm{app}}'(v_{\mathrm{mvt}})
=
\frac{Q_{\mathrm{app}}(z_{\mathrm{rec}})-Q_{\mathrm{app}}(0)}
{z_{\mathrm{rec}}},
\]
and hence
\[
|Q_{\mathrm{app}}(z_{\mathrm{rec}})-Q_{\mathrm{app}}(0)|
=
|Q_{\mathrm{app}}'(v_{\mathrm{mvt}})|\,z_{\mathrm{rec}}.
\]
% lean: CausalSmith.Stat.AnnotationRarearmFrontier.cone_reciprocal_polynomial_gap

\item To bound this derivative, define the polynomial pulled back to the unit interval by
\[
Q_{\mathrm{pull}}(x_{\mathrm{unit}})
:=
Q_{\mathrm{app}}\!\left(\frac B2x_{\mathrm{unit}}+\frac B2\right)
\quad (x_{\mathrm{unit}}\in\mathbb R).
\]
It has degree at most \(L\), and the affine map sends \([-1,1]\) onto \([0,B]\). Thus
\[
\|Q_{\mathrm{pull}}\|_{\infty,[-1,1]}
=
\|Q_{\mathrm{app}}\|_{\infty,[0,B]},
\qquad
Q_{\mathrm{pull}}'(x_{\mathrm{unit}})
=
\frac B2
Q_{\mathrm{app}}'\!\left(\frac B2x_{\mathrm{unit}}+\frac B2\right).
\]
Applying \cref{obj:lem:markov-polynomial} with bound
\(\|Q_{\mathrm{app}}\|_{\infty,[0,B]}\) gives
\[
|Q_{\mathrm{pull}}'(x_{\mathrm{unit}})|
\le
L^2\|Q_{\mathrm{app}}\|_{\infty,[0,B]}
\quad (x_{\mathrm{unit}}\in[-1,1]).
\]
For every \(v\in[0,B]\), the inverse coordinate
\[
x_{\mathrm{unit}}(v):=\frac{2v-B}{B}
\]
belongs to \([-1,1]\). Substituting it and dividing by \(B/2>0\) yields
\[
\|Q_{\mathrm{app}}'\|_{\infty,[0,B]}
\le
\frac{2L^2}{B}\|Q_{\mathrm{app}}\|_{\infty,[0,B]}
\le
\frac{9L^2}{4B}.
\]
% lean: Causalean.Mathlib.Analysis.FinitePolynomialAlternationDuality.markov_derivative_Icc

In particular,
\[
|Q_{\mathrm{app}}'(v_{\mathrm{mvt}})|\,z_{\mathrm{rec}}
\le
\frac{9L^2}{4B}\,z_{\mathrm{rec}}
=
\frac{9}{400c_\epsilon}
\le
\frac{3}{100}
<
\frac14.
\]
This contradicts the chain
\[
\frac14
<
Q_{\mathrm{app}}(0)-Q_{\mathrm{app}}(z_{\mathrm{rec}})
\le
|Q_{\mathrm{app}}(z_{\mathrm{rec}})-Q_{\mathrm{app}}(0)|
=
|Q_{\mathrm{app}}'(v_{\mathrm{mvt}})|\,z_{\mathrm{rec}}.
\]
Therefore
\[
\|g_{\mathrm{rec}}-Q_{\mathrm{app}}\|_{\infty,[0,B]}
\ge\frac18,
\qquad
E_{\mathrm{app}}\ge\frac18.
\]
% lean: CausalSmith.Stat.AnnotationRarearmFrontier.cone_reciprocal_polynomial_gap

\item Apply \cref{obj:lem:finite-polynomial-duality} to the continuous function \(g_{\mathrm{rec}}\) on the nondegenerate interval \([0,B]\). Indexing its nodes from zero, it supplies
\[
0\le v_0<\cdots<v_{L+1}\le B,
\qquad
c_0,\ldots,c_{L+1}\in\mathbb R,
\]
with
\[
\sum_{i=0}^{L+1}|c_i|=1,
\qquad
\sum_{i=0}^{L+1}c_i v_i^h=0
\quad (h=0,\ldots,L),
\qquad
\sum_{i=0}^{L+1}c_i g_{\mathrm{rec}}(v_i)=E_{\mathrm{app}}.
\]
Define
\[
\sigma:=\sum_{i=0}^{L+1}c_i\delta_{v_i}.
\]
Because the nodes are distinct, its total variation and moments satisfy
\[
\|\sigma\|_{\mathrm{TV}}
=
\sum_{i=0}^{L+1}|c_i|=1,
\qquad
\int v^h\,d\sigma(v)
=
\sum_{i=0}^{L+1}c_i v_i^h=0
\quad (h=0,\ldots,L).
\]
Finally,
\[
\int\frac{\alpha}{\alpha+(1-\epsilon)v}\,d\sigma(v)
=
\sum_{i=0}^{L+1}c_i g_{\mathrm{rec}}(v_i)
=
E_{\mathrm{app}}
\ge\frac18,
\]
which establishes the certificate. % lean: CausalSmith.Stat.AnnotationRarearmFrontier.shrinking_cone_dual
\end{enumerate}
\end{proof}
